\section{Discussion of the Extensions}\label{sec:extension_limits}
The selected-subset recovery comparison provides direct evidence of improved functional
reliability on previously problematic inputs. On the selected x64 cohort,
working rate rises from 1.2\% to 80.2\%; x86 also improves, from 14.9\%
to 25.9\%. The pipeline updates therefore recover many useful outputs,
with a substantially stronger improvement on x64. Results on the full
corpus complement this comparison by evaluating the added wrapping and
LLM modes across all 102 original payloads.

The controlled suites help explain what each stage contributes. All six
annotated programs are classified exactly, while the 13 functional
programs check the output of individual programming patterns. Their small
size and simple layouts make them useful diagnostic tests. Their high
success rates are not a consequence of being benign: the original corpus
contains larger and more varied instruction/data layouts, which exercise
additional cases. The recorded native x86 IBP result remains associated with its tested
version; a future corrected implementation requires a separately identified
evaluation.

The two LLM roles remain distinct. Identifying data supplies regions to
the deterministic pipeline; transforming assembly supplies a proposed
instruction listing. On 19 selected original payloads, the former
produces the same final binary as symbolic preparation in 16 cases and
the same 17/19 functional result. For transformation, zero-shot and two-shot
produce similar aggregate outcomes, with no consistent observed advantage
from adding two examples. Input size, construction success, and execution
outcome provide more informative distinctions in this campaign.

The empirical scope can be stated once. Functional preservation is assessed
through the specified execution tests, rather than a formal proof for every
possible execution. Classification accuracy is measured on independently
annotated programs; the unannotated payload comparison measures agreement
and downstream functionality. Static detection rates describe the submitted
raw files and the reported engine responses, including the disclosed
unsupported and failed analyses. Repeated inputs, shared hashes, recorded
version differences, and retries are retained in the artifact so that these
descriptive results can be inspected and reproduced.
